% !TeX encoding = UTF-8
% !TeX program = xelatex
% ============================================================================
% 中国地质大学（武汉）研究生学位论文 LaTeX 模板主文件
% 依据：《中国地质大学（武汉）研究生学位论文写作规范》（2025年5月修改版）
%
% 可选类参数说明：
%   degree = master(学硕，默认) | doctor(博士) | profmaster(全日制专硕)
%            | parttime(非全日制专硕) | eqmaster(同等学力硕士) | eqdoctor(同等学力博士)
%   bibstyle = numerical(顺序编码制，默认) | authoryear(著者-出版年制)
%   fontset  = auto(自动检测，默认) | mac | windows | fandol
%   cover    = true(输出中文封面，默认) | false(学校统一印制封面时关闭)
%   spine    = false(默认) | true(末尾附加竖排书脊页)
%   blind    = false(默认) | true(盲审模式，自动隐藏个人信息)
%   print    = false(电子版，默认) | true(双面印刷模式，中文摘要从奇数页起)
%   frontheader = true(前置页显示页眉，默认) | false(前置页仅显示外侧罗马页码)
% ============================================================================
\documentclass[
  degree=master,
  bibstyle=numerical,
  fontset=auto,
  cover=true,
  spine=false,
  blind=false,
  print=false
]{cugthesis}

% 示例插图使用 TikZ 绘制矢量图（实际写作可直接 \includegraphics 插入 pdf/png/eps）
\usepackage{tikz}
\usetikzlibrary{arrows.meta, positioning, shapes.geometric}

% ============================================================================
% 论文基本元数据配置（所有封面、题名页、声明、页眉均由此自动生成）
% ============================================================================
\cugsetup{
  % 分类号（《中国图书馆分类法》第4版）与 UDC（www.udcc.org）
  clc               = {P631.4},
  udc               = {550.8},
  % 密级：公开论文可填“公开”或留空；涉密论文经批准后标注如“保密★2年”
  secret            = {公开},
  % 学校代码固定为 10491
  schoolcode        = {10491},
  studentid         = {1202310001},
  % 论文题目（规范1.1：一般不宜超过25字，必要时填写 subtitle）
  title             = {基于多尺度特征融合的地球物理数据智能反演方法研究},
  subtitle          = {},
  % 长题目页眉简称（留空则自动使用 title）
  header-title      = {},
  % 英文论文题目
  title-en          = {Research on Intelligent Inversion of Geophysical Data Based on Multi-Scale Feature Fusion},
  subtitle-en       = {},
  % 作者姓名（中英文）
  author            = {张三},
  author-en         = {Zhang San},
  % 学科专业 / 专业学位类别（以教育部2022版目录为准）
  major             = {计算机科学与技术},
  major-en          = {Computer Science and Technology},
  % 英文学位门类（如 Engineering / Science / Computer Technology）
  degree-subject-en = {Engineering},
  % 指导教师（姓名 + 空格 + 职称；若有副导师/企业导师按需填写）
  supervisor        = {李\ 四\quad 教授},
  supervisor-en     = {Prof. Li Si},
  cosupervisor      = {},
  corpsupervisor    = {},
  corpsupervisor-en = {},
  % 培养单位（规范全称，如地球科学学院、计算机学院等）
  department        = {计算机学院},
  % 论文完成时间（规范1.1/3.2/3.4：必须用汉字，不用阿拉伯数字）
  date              = {二○二六年五月},
  % 关键词（3~5个，中文用全角分号“；”分隔，英文用半角分号“; ”分隔）
  keywords          = {地球物理反演；深度学习；多尺度特征融合；物理约束；不确定性量化},
  keywords-en       = {Geophysical Inversion; Deep Learning; Multi-Scale Feature Fusion; Physics Constraints; Uncertainty Quantification},
  % 可选：电子手写签名图片路径及签字日期（纸质提交需亲笔签名，平时留空即可）
  author-sig        = {},
  supervisor-sig    = {},
  sig-date          = {},
}

% ============================================================================
% 正文开始（严格遵循规范第1章规定的 16 个组成部分先后顺序）
% ============================================================================
\begin{document}

% 1. 中文封面  2. 中文题名页  3. 英文题名页
% 4. 原创性声明  5. 导师承诺书  6. 使用授权书
\makefrontpages

% 7. 作者简介（规范1.6、附件6）
% !TeX root = ../main.tex
% 附件6：作者简介（标题格式同各章标题，正文宋体小四、固定行距20磅）
\begin{resume}[作者简介]

\noindent\textbf{一、基本情况}

\noindent 姓名：张三\hspace{0.9em}性别：男\hspace{0.9em}民族：汉\hspace{0.9em}出生年月：2001-05-18\hspace{0.9em}籍贯：湖北省武汉市

\noindent 2019.09—2023.06 \quad 中国地质大学（武汉）计算机学院 \quad 工学学士

\noindent 2023.09—2026.06 \quad 中国地质大学（武汉）计算机学院 \quad 攻读硕士学位

\vspace{6bp}
\noindent\textbf{二、学术论文}

\begin{enumerate}[label=\arabic*., leftmargin=1.8em]
  \item 张三，李四．基于多尺度特征融合的地球物理数据智能反演方法[J]．地球物理学报，2025，68（4）：1234--1246．
  \item Zhang S, Li S, Wang W. Physics-informed multi-scale neural network for seismic impedance inversion[J]. Geophysics, 2025, 90(3): R101--R115.（已录用待发表）
\end{enumerate}

\vspace{6bp}
\noindent\textbf{三、获奖、专利情况}

\begin{enumerate}[label=\arabic*., leftmargin=1.8em]
  \item 张三，李四．一种基于物理约束的多尺度地震反演系统及方法：中国，\allowbreak CN202410123456.7[P]．2025-02-18．
  \item 2024 年中国地质大学（武汉）研究生国家奖学金．
\end{enumerate}

\vspace{6bp}
\noindent\textbf{四、研究项目}

\begin{enumerate}[label=\arabic*., leftmargin=1.8em]
  \item 国家自然科学基金面上项目：面向复杂构造区的深度学习地球物理反演关键技术研究，2023.01—2026.12，项目编号：42274000，主要参与人．
\end{enumerate}

\end{resume}


% 8. 学位论文答辩委员会名单（规范1.7、附件7，置于中文摘要之前）
% !TeX root = ../main.tex
% 附件7：学位论文答辩委员会名单（置于中文摘要之前）
\begin{committee}
  \committeemember{主\quad 席}{王五}{教授}{武汉大学}
  \committeemember{委\quad 员}{赵六}{教授}{华中科技大学}
  \committeemember{}{钱七}{教授}{中国地质大学（武汉）}
  \committeemember{}{孙八}{副教授}{中国地质大学（武汉）}
  \committeemember{}{周九}{副教授}{中国地质大学（武汉）}
  \committeemember{秘\quad 书}{吴十}{讲师}{中国地质大学（武汉）}
\end{committee}


% 9. 中文摘要（自此起启用罗马数字页码 I, II, III...，规范3.10自此双面印刷）
% !TeX root = ../main.tex
% 第1.8节、第3.6节：中文摘要（篇幅以一页为宜，不出现图表、化学式和非公知术语）
\begin{abstractzh}
地球物理反演是地下介质结构成像与矿产资源勘探的核心技术手段。然而，传统基于最优化迭代的确定性反演方法通常面临计算开销大、对初始模型高度敏感以及强非线性条件下容易陷入局部极小值等瓶颈问题。近年来，数据驱动的深度学习方法在地球物理反演领域展现出显著优势，但在实际复杂地质场景中仍存在多尺度构造特征提取不足、泛化能力受限以及违背地球物理正演控制方程等突出挑战。

针对上述问题，本文围绕复杂地下介质的高精度智能反演展开系统研究，主要研究工作与创新性成果包括以下几个方面：

第一，提出了一种融合残差注意力机制的多尺度编解码网络结构。通过在编码端引入跨尺度空洞卷积金字塔与通道空间联合注意力模块，实现了对地下大尺度背景速度场与小尺度薄互层异常体的协同表征，有效缓解了深层网络下采样过程中的高频细节丢失问题。

第二，构建了嵌入正演物理算子的双驱动联合约束反演框架。将一维褶积与二维波动方程正演过程封装为可微物理层并嵌入训练损失函数，在不显著增加推理耗时的前提下，强制反演结果满足观测数据的动力学与运动学一致性，从而大幅提升了模型在少样本及跨工区迁移条件下的泛化能力。

第三，设计了基于蒙特卡洛随机失活与集成学习相结合的不确定性量化策略，在输出高分辨率地下物性参数估计的同时，定量给出各空间位置反演结果的置信区间与后验方差分布，为后续地质解释与钻井部署提供可靠的决策依据。

在多个合成基准数据集及某实际三维工区观测资料上的对比实验表明，与传统全波形反演方法及现有主流深度学习反演模型相比，本文方法在反演精度、抗噪鲁棒性以及计算效率方面均取得了显著提升，反演剖面的决定系数提高约 6.4\%，单道推理速度提升两个数量级以上，具有良好的工程应用价值。
\end{abstractzh}


% 10. Abstract
% !TeX root = ../main.tex
% 第1.9节、第3.6节：Abstract（不少于300个英文实词，与中文摘要相对应）
\begin{abstracten}
Geophysical inversion is a fundamental technique for subsurface imaging and mineral exploration. However, conventional iterative optimization methods often suffer from high computational cost, sensitivity to initial models, and local minima under strong nonlinearity. Although deep learning has shown great promise in geophysical inversion, practical deployment in complex geological settings remains limited by insufficient multi-scale feature extraction, weak cross-region generalization, and potential violations of governing physical equations.

To address these challenges, this thesis investigates high-precision intelligent inversion for complex subsurface media. The main contributions are summarized as follows:

First, a multi-scale encoder-decoder network with residual attention is proposed. By combining cross-scale dilated convolution pyramids with channel-spatial attention modules, the network captures both large-scale background trends and fine-scale thin-bed anomalies while mitigating high-frequency detail loss.

Second, a physics-constrained dual-driven inversion framework is constructed by embedding differentiable forward operators into the training loss. Without noticeably increasing inference latency, this enforces kinematic and dynamic consistency between predictions and observations, improving generalization under scarce labels and cross-survey transfer.

Third, an uncertainty quantification strategy integrating Monte Carlo dropout and deep ensembles is developed to estimate spatial confidence intervals and posterior variance maps alongside subsurface property predictions.

Experiments on synthetic benchmarks and a real 3D field survey show that the proposed method outperforms conventional full-waveform inversion and existing deep learning baselines in accuracy, noise robustness, and speed, improving $R^2$ by approximately 6.4\% while accelerating inference by over two orders of magnitude.
\end{abstracten}


% 11. 目录（规范1.10、3.7、附件8）
\cugtableofcontents

% 12. 图和表清单（规范1.11、附件9；非必要项，如不需要直接注释以下两行）
\cuglistoffigures
\cuglistoftables

% 13. 正文（自第一章起阿拉伯数字页码 1, 2, 3...，规范2.4页眉及外侧页码）
\cugmainmatter
% !TeX root = ../main.tex
\chapter{绪论}

\section{研究背景及意义}

地球物理反演通过地表或井中观测到的地震波、重磁、电磁等地球物理场响应，重建地下介质的物性参数空间分布，是深部资源勘探、灾害预警和城市地下空间探测的核心科学问题\cite{tarantola2005inverse, zhao1995mechanics}。随着勘探目标逐渐向深层、隐蔽油气藏以及复杂构造区转移，地下介质呈现出显著的非均匀性、各向异性与多尺度嵌套特征，对反演成像的分辨率与可靠性提出了更高要求\footnote{本模板脚注按规范第 1.12 节采用当页脚注方式，每页自动从带圈数字 \ding{172} 开始重新编号，采用小五号宋体排版。}。

传统反演方法以最小二乘优化与伴随状态法为代表，在理论体系上较为完备，但在面对强非线性、观测孔径受限及含噪数据时存在严重的不适定性\cite{gbt7714, wang2023seismic}。另一方面，伴随深度学习在图像重建与科学计算中的兴起，融合领域先验知识与数据驱动模型的智能反演方法成为当前地球物理与人工智能交叉领域的研究热点\cite{raissi2019physics, zhang2024inversion}。

\section{国内外研究现状}

\subsection{传统确定性地球物理反演方法}

经典线性化反演与全波形反演（Full-Waveform Inversion, FWI）通过最小化合成记录与实测数据之间的残差泛函迭代更新地下模型参数\cite{tarantola2005inverse}。为克服反演多解性，研究者相继引入了 Tikhonov 正则化、全变差（Total Variation）约束以及多尺度频带延拓策略\cite{desmarais1992carbon, yang2021fwi}。尽管上述方法在低频信息充足时能够获得高精度成像结果，但其庞大的正演模拟开销与严苛的初始模型要求限制了实时成像与大规模三维应用。

\subsection{基于深度学习的智能反演方法}

近年来，卷积神经网络、U-Net 及其变体被广泛应用于地震波阻抗反演与速度建模任务\cite{wu2019inversionnet, li2022dl}。然而，纯数据驱动方法在训练样本分布与目标工区存在差异时容易出现预测漂移。为提升物理一致性，物理信息神经网络（Physics-Informed Neural Networks, PINN）将偏微分方程残差直接引入损失函数\cite{raissi2019physics}，为高泛化智能反演提供了新思路。

\section{本文主要研究内容与结构安排}

围绕上述关键科学问题，本文共分为四章，具体组织结构如下：
\begin{enumerate}[label=（\arabic*）]
  \item 第一章为绪论，阐述研究背景、国内外研究现状及本文的总体框架；
  \item 第二章为研究方法与理论分析，建立多尺度特征融合网络与物理约束联合反演模型，并演示公式、图、三线表与跨页续表的规范排法；
  \item 第三章为实验结果与讨论，结合合成数据与实际工区资料开展定量评价与消融分析；
  \item 第四章为结论与展望，总结全文主要创新点，指出当前研究局限并展望未来工作方向。
\end{enumerate}

% !TeX root = ../main.tex
\chapter{研究方法与理论分析}

\section{地球物理反演的数学建模}

地球物理正演过程可以抽象为从模型空间 $\mathcal{M}$ 到数据空间 $\mathcal{D}$ 的非线性映射算子 $\mathcal{F}$，其含噪观测模型可表示为
\begin{equation}
  \mathbf{d}_{\mathrm{obs}} = \mathcal{F}(\mathbf{m}) + \boldsymbol{\varepsilon},
  \label{eq:forward}
\end{equation}
其中，$\mathbf{d}_{\mathrm{obs}} \in \mathbb{R}^{N_t \times N_r}$ 表示实测地球物理数据矩阵，$\mathbf{m} \in \mathbb{R}^{H \times W}$ 表示待反演地下介质物性参数场，$\boldsymbol{\varepsilon}$ 为加性观测噪声。在物理量表达中，波速单位记为 \si{\meter\per\second}，密度单位记为 \si{\kilogram\per\cubic\meter}，频率单位记为 \si{\hertz}，严格遵循 GB 3100～3102—1993 量和单位国家标准（量符号用斜体，单位符号用正体，数值与单位之间保留适当间隙，如 \SI{3500}{\meter\per\second}、\SI{2.65e3}{\kilogram\per\cubic\meter}）。

为同时利用有标签数据先验与物理控制方程约束，本文构造如下复合目标函数：
\begin{equation}
  \mathcal{L}(\theta) = \frac{1}{N} \sum_{i=1}^{N} \left\| \hat{\mathbf{m}}_i - \mathbf{m}_i \right\|_2^2
  + \lambda_{\mathrm{phy}} \left\| \mathcal{F}(\hat{\mathbf{m}}_i) - \mathbf{d}_i \right\|_2^2
  + \lambda_{\mathrm{reg}} \left\| \nabla \hat{\mathbf{m}}_i \right\|_1,
  \label{eq:loss}
\end{equation}
式中，$\theta$ 为神经网络待优化权重参数，$\hat{\mathbf{m}}_i = \mathcal{G}_\theta(\mathbf{d}_i)$ 为网络输出的反演结果，$\lambda_{\mathrm{phy}}$ 与 $\lambda_{\mathrm{reg}}$ 分别为物理一致性约束系数与梯度稀疏正则化系数。

\section{多尺度特征融合反演网络}

如图~\ref{fig:arch} 所示，本文提出的多尺度物理约束反演框架由多尺度编码器、跨层残差注意力桥接模块、渐进式上采样解码器以及可微正演物理层四部分组成。图序号与图名置于图下方，采用五号宋体居中排版，序号与图题之间空一个汉字符宽度。

\begin{figure}[htbp]
  \centering
  \begin{tikzpicture}[
    node distance=0.85cm and 0.9cm,
    box/.style={draw, rounded corners=2pt, minimum width=2.4cm, minimum height=0.95cm, align=center, font=\songti\zihao{5}, fill=gray!6},
    phy/.style={draw, dashed, rounded corners=2pt, minimum width=2.6cm, minimum height=0.95cm, align=center, font=\songti\zihao{5}, fill=gray!14},
    arr/.style={-{Stealth[length=2.2mm]}, thick}
  ]
    \node[box] (in) {观测地震数据\\$\mathbf{d}_{\mathrm{obs}}$};
    \node[box, right=of in] (enc) {多尺度编码器\\（空洞金字塔）};
    \node[box, right=of enc] (att) {残差注意力\\特征融合};
    \node[box, right=of att] (dec) {渐进式解码器\\$\hat{\mathbf{m}}$};
    \node[phy, below=1.0cm of dec] (fwd) {可微正演物理算子\\$\mathcal{F}(\hat{\mathbf{m}})$};
    \node[phy, below=1.0cm of enc] (loss) {数据--物理联合损失\\$\mathcal{L}(\theta)$};

    \draw[arr] (in) -- (enc);
    \draw[arr] (enc) -- (att);
    \draw[arr] (att) -- (dec);
    \draw[arr] (dec) -- (fwd);
    \draw[arr] (fwd) -- (loss);
    \draw[arr] (in.south) |- (loss.west);
    \draw[arr, dashed] (loss.north) -- node[right, font=\zihao{-5}] {反向传播} (enc.south);
  \end{tikzpicture}
  \caption{物理约束多尺度特征融合反演网络总体架构示意图}
  \label{fig:arch}
\end{figure}

\section{关键参数设置与跨页续表示例}

表~\ref{tab:hyperparams} 给出了模型训练阶段的核心超参数配置。依据规范第 2.5.2 节与表 3.8 要求，表采用国际通用三线表格式，表序号与表题目置于表上方，采用五号宋体居中，序号与题名之间空一个汉字符宽度。

\begin{table}[htbp]
  \centering
  \caption{模型训练超参数与正演模拟关键参数设置}
  \label{tab:hyperparams}
  \begin{tabular}{cccc}
    \toprule
    参数类别 & 参数名称 & 设定值 & 单位 / 说明 \\
    \midrule
    正演参数 & 震源主频 $f_0$ & 25.0 & \si{\hertz} \\
    正演参数 & 空间网格步长 $\Delta x, \Delta z$ & 10.0 & \si{\meter} \\
    正演参数 & 时间采样间隔 $\Delta t$ & 1.0 & \si{\milli\second} \\
    训练参数 & 初始学习率 $\eta_0$ & $1.0\times 10^{-3}$ & 余弦退火衰减 \\
    训练参数 & 物理约束权重 $\lambda_{\mathrm{phy}}$ & 0.25 & 无量纲 \\
    训练参数 & 正则化权重 $\lambda_{\mathrm{reg}}$ & 0.05 & 无量纲 \\
    \bottomrule
  \end{tabular}
\end{table}

当表格内容较长需要跨页接排时，规范第 2.5.2 节要求在续页重复表序号、表题（可省略）加“（续）”，并重复表头。表~\ref{tab:longexample} 给出了基于 \texttt{longtable} 宏包的标准续表实现示范。

\begin{longtable}{ccccc}
  \caption{不同岩性层位典型地球物理参数统计表}\label{tab:longexample}\\
  \toprule
  编号 & 岩性名称 & 纵波速度 / (\si{\meter\per\second}) & 密度 / (\si{\gram\per\cubic\centi\meter}) & 孔隙度 / \% \\
  \midrule
  \endfirsthead

  \multicolumn{5}{c}{\songti\zihao{5} 表\thetable\hspace{1\ccwd}不同岩性层位典型地球物理参数统计表（续）}\\[4pt]
  \toprule
  编号 & 岩性名称 & 纵波速度 / (\si{\meter\per\second}) & 密度 / (\si{\gram\per\cubic\centi\meter}) & 孔隙度 / \% \\
  \midrule
  \endhead

  \midrule
  \multicolumn{5}{r}{\songti\zihao{-5}（续下页）}\\
  \endfoot

  \bottomrule
  \endlastfoot

  1  & 第四系松散砂土 & 1650 & 1.85 & 32.4 \\
  2  & 泥质粉砂岩     & 2480 & 2.18 & 21.6 \\
  3  & 细粒石英砂岩   & 3120 & 2.34 & 16.2 \\
  4  & 钙质胶结砂岩   & 3680 & 2.49 & 11.5 \\
  5  & 致密页岩       & 3350 & 2.52 &  6.8 \\
  6  & 生物碎屑灰岩   & 4850 & 2.66 &  4.9 \\
  7  & 白云质灰岩     & 5420 & 2.74 &  3.7 \\
  8  & 结晶白云岩     & 5890 & 2.81 &  2.5 \\
  9  & 火山凝灰岩     & 4120 & 2.55 &  8.3 \\
  10 & 安山玄武岩     & 5640 & 2.86 &  1.9 \\
  11 & 花岗闪长岩     & 5980 & 2.72 &  1.1 \\
  12 & 变质片麻岩     & 6150 & 2.79 &  0.8 \\
\end{longtable}

% !TeX root = ../main.tex
\chapter{实验结果与讨论}

\section{实验设置与评价指标}

为全面评估所提方法的有效性与泛化能力，本文选取决定系数 $R^2$、均方根误差（RMSE）及结构相似性指数（SSIM）作为定量评价指标。其中，均方根误差定义为
\begin{equation}
  \mathrm{RMSE} = \sqrt{\frac{1}{HW} \sum_{p=1}^{H} \sum_{q=1}^{W} \left( \hat{m}_{p,q} - m_{p,q} \right)^2}.
  \label{eq:rmse}
\end{equation}

\section{对比实验与消融分析}

\subsection{不同反演方法精度对比}

将本文提出的多尺度物理约束网络（MS-PINN）与传统全波形反演（FWI）、标准 U-Net 及 InversionNet 在含噪测试集上进行对比，定量指标汇总见表~\ref{tab:compare}。

\begin{table}[htbp]
  \centering
  \caption{不同反演方法在测试集上的综合性能对比}
  \label{tab:compare}
  \begin{tabular}{lcccc}
    \toprule
    反演方法 & $R^2$ & RMSE / (\si{\meter\per\second}) & SSIM & 单样本耗时 / \si{\second} \\
    \midrule
    传统 FWI\cite{tarantola2005inverse}   & 0.861 & 214.5 & 0.842 & 128.40 \\
    标准 U-Net\cite{li2022dl}              & 0.884 & 186.2 & 0.875 &   0.04 \\
    InversionNet\cite{wu2019inversionnet}  & 0.902 & 162.8 & 0.896 &   0.05 \\
    本文方法（MS-PINN）                    & \textbf{0.960} & \textbf{103.4} & \textbf{0.952} & \textbf{0.06} \\
    \bottomrule
  \end{tabular}
\end{table}

\subsection{抗噪鲁棒性分析}

图~\ref{fig:noise} 展示了在不同信噪比（SNR）条件下各方法的收敛特性与误差变化趋势。由图~\ref{fig:noise}(a) 与图~\ref{fig:noise}(b) 可知，得益于物理约束项对高频随机干扰的抑制作用，当信噪比降低至 \SI{10}{\decibel} 时，本文方法仍能维持较高的结构保真度。

\begin{figure}[htbp]
  \centering
  \begin{subfigure}[b]{0.46\textwidth}
    \centering
    \begin{tikzpicture}[scale=0.82]
      \draw[->] (0,0) -- (5.2,0) node[right, font=\zihao{5}] {Epoch};
      \draw[->] (0,0) -- (0,3.4) node[above, font=\zihao{5}] {Loss};
      \draw[thick, dashed] plot[smooth] coordinates {(0.3,3.0) (1.2,2.1) (2.4,1.5) (3.6,1.2) (4.8,1.05)};
      \draw[thick] plot[smooth] coordinates {(0.3,2.9) (1.2,1.6) (2.4,0.95) (3.6,0.65) (4.8,0.52)};
      \node[font=\zihao{-5}] at (3.6,1.55) {U-Net};
      \node[font=\zihao{-5}] at (3.6,0.32) {本文方法};
    \end{tikzpicture}
    \caption{训练损失收敛曲线}
  \end{subfigure}
  \hfill
  \begin{subfigure}[b]{0.46\textwidth}
    \centering
    \begin{tikzpicture}[scale=0.82]
      \draw[->] (0,0) -- (5.2,0) node[right, font=\zihao{5}] {SNR / dB};
      \draw[->] (0,0) -- (0,3.4) node[above, font=\zihao{5}] {$R^2$};
      \draw[thick, dashed] plot[smooth] coordinates {(0.5,1.2) (1.6,1.8) (2.8,2.25) (4.0,2.5) (4.8,2.65)};
      \draw[thick] plot[smooth] coordinates {(0.5,1.95) (1.6,2.45) (2.8,2.85) (4.0,3.05) (4.8,3.15)};
      \node[font=\zihao{-5}] at (3.7,2.15) {U-Net};
      \node[font=\zihao{-5}] at (3.5,3.22) {本文方法};
    \end{tikzpicture}
    \caption{不同信噪比下的反演精度}
  \end{subfigure}
  \caption{模型收敛特性与抗噪鲁棒性对比分析}
  \label{fig:noise}
\end{figure}

% !TeX root = ../main.tex
\chapter{结论与展望}

\section{主要结论}

本文针对复杂地质构造背景下地球物理反演精度不足与泛化能力受限的问题，开展了多尺度特征融合网络构建、物理算子嵌入约束以及不确定性量化三个方面的研究，得出如下主要结论：

\begin{enumerate}[label=（\arabic*）]
  \item 引入跨尺度空洞卷积金字塔与残差注意力机制后，反演网络对薄互层与断裂带等细微地质异常的刻画能力明显增强，有效克服了深层特征提取中的高频衰减现象。
  \item 将可微正演物理算子融入端到端训练闭环，在保持毫秒级推理效率的同时显著降低了反演解的物理失配度，增强了模型在低信噪比与跨工区场景下的泛化稳定性。
  \item 结合蒙特卡洛随机失活与模型集成构建的后验方差估计能够准确指示低置信度反演区域，为地球物理定量解释提供了直观的可靠性度量。
\end{enumerate}

\section{研究局限与未来展望}

尽管本文方法在多个实验场景中取得了理想效果，但仍存在若干值得深入探索的方向：其一，当前可微物理层主要面向二维声波近似，未来可进一步扩展至三维弹性波及各向异性介质情形；其二，如何在多模态地球物理资料（如地震--重磁--电磁联合观测）下构建统一的基础大模型，将是下一步研究的重点内容。


% 14. 参考文献（规范1.13、2.6、3.9、附件11，GB/T 7714）
\cugbackmatter
\cugbibliography{bib/references}

% 15. 致谢（规范1.14、3.9）
% !TeX root = ../main.tex
% 第1.14节、第3.9节：致谢（标题同各章标题，正文宋体小四、固定行距20磅）
\begin{acknowledgements}
时光荏苒，白驹过隙。三年的硕士研究生求学生涯即将画上圆满的句号，在南望山麓、未来城畔的中国地质大学（武汉）度过的这段时光，将成为我人生中最珍贵的记忆。

首先，衷心感谢我的指导教师李四教授。从论文选题、理论推演、实验设计到论文撰写与反复修改，导师始终给予我悉心的指导与无私的关怀。导师严谨治学的科学态度、深厚渊博的学术造诣以及精益求精的工作作风使我受益终身。

感谢学院各位老师在课程学习、开题答辩、中期考核及论文评审中提出的宝贵意见与建设性建议；感谢课题组全体师兄师姐、同门及师弟师妹在科研讨论与日常生活中给予的真诚帮助与支持。

感谢国家自然科学基金项目（编号：42274000）对本研究工作的资助。

最后，深深感谢我的父母及家人多年来默默的理解、鼓励与付出，是你们温暖的陪伴让我得以心无旁骛地完成学业。谨向所有关心和帮助过我的师长、亲友致以最诚挚的谢意！
\end{acknowledgements}


% 16. 附录（规范1.15、3.9；如无附录可注释）
\cugappendix
% !TeX root = ../main.tex
% 第1.15节、第3.9节：附录（标题同各章标题，正文宋体小四、固定行距20磅）
\chapter{主要符号说明与补充推导}

本附录给出正文中涉及的主要数学符号定义（表~\ref{tab:symbols}），以方便读者查阅。

\begin{table}[htbp]
  \centering
  \caption{论文主要数学符号及物理含义说明}
  \label{tab:symbols}
  \begin{tabular}{ccc}
    \toprule
    符号 & 含义 & 单位 \\
    \midrule
    $\mathbf{d}_{\mathrm{obs}}$ & 地表观测地球物理数据矩阵 & 无量纲 / \si{\meter\per\second} \\
    $\mathbf{m}, \hat{\mathbf{m}}$ & 真实与反演地下物性参数模型 & \si{\meter\per\second} \\
    $\mathcal{F}(\cdot)$ & 地球物理正演映射算子 & — \\
    $\mathcal{G}_\theta(\cdot)$ & 参数为 $\theta$ 的深度学习反演网络 & — \\
    $\lambda_{\mathrm{phy}}, \lambda_{\mathrm{reg}}$ & 物理一致性与正则化权重系数 & 无量纲 \\
    \bottomrule
  \end{tabular}
\end{table}

在弱扰动线性化假设下，正演残差关于模型参数的一阶 Fr\'echet 导数矩阵（即雅可比矩阵）记为 $\mathbf{J} = \partial \mathcal{F}(\mathbf{m}) / \partial \mathbf{m}$，由此可得高斯--牛顿型迭代更新增量：
\begin{equation}
  \Delta \mathbf{m} = -\left( \mathbf{J}^{\mathsf{T}}\mathbf{J} + \alpha \mathbf{L}^{\mathsf{T}}\mathbf{L} \right)^{-1}
  \mathbf{J}^{\mathsf{T}}\left( \mathcal{F}(\mathbf{m}) - \mathbf{d}_{\mathrm{obs}} \right).
  \label{eq:gn}
\end{equation}


\end{document}
